%laden der Präambel mit Latexbefehlen/-klassen
% !TeX encoding = UTF-8
% !TeX spellcheck = de_DE

%Dokumentklasse, Spracheinstellung, Schriften
\documentclass[a0paper,portrait]{baposter}
\usepackage[utf8]{inputenc}
\usepackage{arev}
\usepackage[T1]{fontenc}

%Mathe, Symbole, Einheitendarstellung, Chemie
\usepackage{amsmath}
\usepackage{amsxtra}
\usepackage{eurosym}
\usepackage{siunitx}  
\sisetup{locale=DE}
\usepackage[version=4]{mhchem}

%Typographie
\usepackage[auto]{microtype}

%Einbindung von Bildern, Tabellen, pdf-Seiten, Quellcode
\usepackage{booktabs}
\usepackage{multirow}
\usepackage{paralist}
\usepackage{listings}

%Darstellung der Literaturangaben
\usepackage[
backend=biber,
style=science,
citestyle=numeric-comp,
sorting=none,
maxbibnames=1,
firstinits=true
]{biblatex}

%Fuer affiliations
\usepackage{authblk}
\bibliography{literatur/refs}
\setlength\bibitemsep{2.5pt}

%Fuer workflow flowchart
\usepackage{tikz}
\usetikzlibrary{shapes.geometric, shapes.arrows, arrows}
\usetikzlibrary{positioning}

%Fuer Bild+Unterschriften
\usepackage{graphicx}
%\usepackage{capt-of}
\usepackage[labelformat=empty]{caption}

%Text farblich hervorheben
%\usepackage[usenames,dvipsnames]{xcolor}

%Mathe
\usepackage{bm}
%Definition der Farben
\definecolor{standardfontcolor}{RGB}{0,0,0} 
\definecolor{bordercol}{RGB}{0,200,200}
\definecolor{headercol1}{RGB}{255,255,255}
\definecolor{headercol2}{RGB}{0,200,200}
\definecolor{headerfontcol}{RGB}{0,0,0}
\definecolor{boxcolor}{RGB}{255,255,255}


\begin{document}
	\typeout{Poster rendering started}
	
\background{
	\begin{tikzpicture}[remember picture,overlay]%
	\draw (current page.north west)+(-2em,2em) node[anchor=north west]
	{\includegraphics[height=1.1\textheight]{figures/background}};
	\end{tikzpicture}
}

\color{standardfontcolor}

\begin{poster}{
		grid=false,
		columns=3,
		%colspacing=length
		headerheight=0.125\textheight,
		eyecatcher=false, 
		borderColor=bordercol,
		headerColorOne=headercol1,
		headerColorTwo=headercol2,
		headerFontColor=headerfontcol,
		% Only simple background color used, no shading, so boxColorTwo isn't necessary
		boxColorOne=boxcolor,
		headershape=roundedright,
		headerfont=\sffamily\bfseries\Large,
		textborder=rectangle,
		headerborder=open,
		boxshade=plain,
		background=none
%		background=user
	}
	%%% Eye Cacther %%%%%%%%%%%%%%%%%%%%%%%%%%%%%%%%%%%%%%%%%%%%%%%%%%%%%%%%%%%%%%%
	{
		Eye Catcher, empty if option eyecatcher=false - unused
	}
%----------------------------------------------------------------------------------------
%	TITLE AND AUTHOR NAME
%----------------------------------------------------------------------------------------
%
{
	\textsf %Sans Serif
	{Periodic Coupled Cluster Theory and \\Applications: Interface to FHI-aims and VASP}
}
{\sf\vspace{0.5em}\\
	Evgeny Moerman\textsuperscript{1}, Andreas Irmler\textsuperscript{2}, Felix Hummel\textsuperscript{2}, 
	Andreas Grüneis\textsuperscript{2}, Matthias Scheffler\textsuperscript{1}
	\vspace{0.1em}\\
	\small{\textsuperscript{1} The NOMAD Laboratory at the FHI-MPG, Berlin
	\vspace{0.2em}\\
	\textsuperscript{2} Institute for Theoretical Physics, TU Wien, Vienna}
}
% University/lab logo
{\includegraphics[height=3cm]{bilder/logo/NOMAD_Logo_srgb_print_wlow.jpg}} 
%
%
\headerbox{Abstract}{name=einleitung,column=0,row=0,span=3}{
%Vorlage für die Erstellung eines Posters. Es wird die Klasse baposter verwendet. In der entsprechenden Dokumentation %finden sich Hinweise zur Gestaltung. Die Vorlage bietet ein einfaches Layout, dass zur Bearbeitung der Aufgabenstellung %ausreichend sein sollte. Die Einbindung verschiedener Elemente ist beispielhaft gezeigt.
%
%Hier könnte die Einleitung stehen, also Motivation, Ziel der Arbeit, Aufgabenstellung, etc.
%
%Literaturverweise, ein Buch \cite{epple2009,gruhn1976} und ein Artikel \cite{Klaucke2020,Hofmann2018}

Coupled cluster (CC) theory achieves a high level of accuracy for a wide range of properties in atoms and molecules at affordable computational cost.
	However, only few codes exist for CC calculations of periodic systems. The \texttt{Cc4s} code~\cite{cc4s}
allows for periodic CCSD(T)-level calculations and it has been interfaced to
VASP\cite{kresse1996efficiency} and, recently, to FHI-aims\cite{moerman2022interface, fhiaims}.
Together these interfaces, which cover atom-centered basis and plane-wave codes, can be generalized to 
interface most other ab-initio software packages to \texttt{Cc4s}. This enables a wide range of applications ranging from molecular systems in the gas phase to solids and surfaces.
}

%\headerbox{General Workflow}{name=Workflow,column=0,below=einleitung,span=2}{
%	%\begin{center}
%		\includegraphics[scale=0.5875]{bilder/CC_aims_flowchart.pdf}		
%	%\end{center}
%	\begin{tikzpicture}
%        \matrix(m) [] {
%            \node {\includegraphics[width=4cm]{bilder/FHI-aims-logo.png}}; \\
%            \node {\includegraphics[width=4cm]{bilder/vasp_logo.png}}; \\
%    };
%    %\node[above right = of a] (c) {c};
%    \node [inner sep=1.0mm, fill=red!50, right = of m, style={single arrow, draw=none}, align=center] %(arrow1) 
%    {perform SCF and \\ write ab-initio \\ quantities  to files \\ via interface};
%    
%    \node [inner sep=0pt, right = of arrow1] (filesimg) %{\includegraphics[width=2cm]{bilder/filesicon.jpg}};
%    \node[below= \belowcaptionskip of filesimg,text width=\linewidth]
%    {\captionof{figure}{\texttt{EigenEnergies.dat} \\ \texttt{CoulombVertex.dat} \\ %\texttt{FockOperator.dat} \\ $\dots$} };
%    
%    \node [inner sep=1.0mm, fill=red!50, right = of filesimg, style={single arrow, draw=none}, %align=center] (arrow2) 
%    {perform \\ correlated \\ calculation \\ (MP2, CCSD, ...)};
%    
%    \node [right = of arrow2]{\includegraphics[width=4cm]{bilder/cc4s_logo_cropped.png}};
%\end{tikzpicture}
%}

\headerbox{EOM-CCSD for trans-[C\textsubscript{2}H\textsubscript{2}]\textsubscript{n}}{name=Bandstructure,column=0,below=einleitung,span=1}{
    %\begin{center}
%		\includegraphics[scale=0.5]{bilder/bandstructure.pdf}		
%	\end{center}

\begin{minipage}{0.7\linewidth}
\includegraphics[scale=0.45]{bilder/bandstructure.pdf}
\end{minipage}
\hspace{0.4cm}
\begin{minipage}{0.2\linewidth}
  \begin{center}
  \includegraphics[scale=0.04, angle=90]{bilder/polyacetylene_centered_long_cropped_2nd.png}
  \end{center}
\end{minipage}

    \begin{center}
     Bandstructure obtained via EA- and IP-EOM-CCSD\cite{stanton1993equation} using 
     NAO-VCC-2Z basis set
     \end{center}

%     \begin{center}
%	     \includegraphics[scale=0.04]{bilder/polyacetylene_centered_long_cropped.png}
%     \end{center}
}
\headerbox{CCSD for H\textsubscript{2}O@C\textsubscript{6}H\textsubscript{6}}{name=H2OGraphene,column=0,below=Bandstructure,span=1}{
\begin{minipage}{0.3\linewidth}
  \begin{center}
    \includegraphics[scale=0.03]{bilder/geometry_small_cropped.png}
  \end{center}
\end{minipage}
\hspace{0.5cm}
\begin{minipage}{0.5\linewidth}
\begin{tabular}{ c|c } 
 Method & E\textsubscript{ad}(kcal/mol)\\\hline
 CCSD/CBS(45) & -2.4  \\
 PBE-D3\cite{ajala2019assessment} & -3.2 \\
 B3LYP-D3\cite{ajala2019assessment} & -3.4 \\
 RPA\cite{ajala2019assessment} & -2.5 \\
 DMC\cite{ajala2019assessment} & -2.9
\end{tabular}
\end{minipage}
}



\headerbox{CCSD(T) for C\textsubscript{2}H\textsubscript{4}@Si(cd)}{name=C2H4Si,column=0,below=H2OGraphene,span=1}{
\begin{minipage}{0.2\linewidth}
\includegraphics[scale=0.2]{bilder/poscar_pre.png}
\end{minipage}
\hspace{2.2cm}
\begin{minipage}{0.5\linewidth}
70 atoms with $N_{\rm o}=110$. CCSD and (T) calculations with $N_{\rm v}=1650$
and $N_{\rm v}=550$, respectively\cite{tonner}. \\

\begin{tabular}{ c|c } 
 Method & E\textsubscript{ad}(kcal/mol)\\
	\hline
 PBE & -12.2  \\ 
 PBE-D3 & -17.1  \\ 
 HF & 1.6 \\
 MP2 & -20.0 \\
 CCSD & -11.4 \\
 CCSD(T) & -9.9
\end{tabular}
\end{minipage}


}

\headerbox{CCSD(T):RPA for H\textsubscript{2}O@[C\textsubscript{3}N\textsubscript{4}]}{name=H2OC3N4,column=0,below=C2H4Si,span=1}{
138 atoms with $N_{\rm o}=292$. Embedded CCSD(T) calculations with
up to $N_{\rm o}=40$ and $N_{\rm v}=800$\cite{schafer2021}.\\
\begin{minipage}{0.18\linewidth}
\vspace{0.9cm}
 \includegraphics[scale=0.18]{./bilder/h2o_gc3n4_topview.png}\\
\bigskip
\includegraphics[scale=0.054]{./bilder/h2o_gc3n4_sideview.png}
\end{minipage}
\hspace{1.8cm}
\begin{minipage}{0.5\linewidth}
\begin{tabular}{ c|c } 
	{\small  Method } & {\small  E\textsubscript{ad} (kcal/mol) }\\
	\hline
 PBE & -8.1  \\ 
 PBE-D3 & -12.2 \\ 
 PBE-TS & -12.7 \\ 
 PBE0-MBD & -12.4 \\ 
 HF & -4.8 \\
	{\small  MP2:dMP2 } & -12.4 \\
	{\small  CCSD:RPA } & -10.7 \\
	{\small  CCSD(T):RPA } & -11.9
\end{tabular}
\end{minipage}


}



%\headerbox{The Coulomb vertex}{name=CoulombVertex,column=0, below=Workflow, span=2}{
%The common memory bottleneck of wave-function based methods is the calculation
%and saving of the Coulomb integrals $V_{pq}^{rs}$.\\
%
%\textcolor[rgb]{0.05, 0.6, 0.35}{Solution: } The Coulomb vertex $\Gamma_{p}^{rG}$, a rank-3 tensor %approximation
%of the $V_{pq}^{rs}$-tensor:
%\begin{equation}
%    V_{pq}^{rs} = \sum_{G} \Gamma_{p}^{rG*}\Gamma_{q}^{sG}
%\end{equation}
%
%Due to the reduced dimension, 
%the memory scaling of such methods is lower, allowing for bigger systems to be calculated.
%}
\headerbox{Atomic orbitals - FHI-aims}{name=fhiaims,column=1,below=einleitung,span=1}{
FHI-aims\cite{fhiaims} is an all-electron electronic structure
package and employs numerical atomic orbitals (NAOs) for both
molecules and periodic systems,
interfaced to \texttt{Cc4s} via the CC-aims interface\cite{moerman2022interface,ccaims}. \\

Interface requirements:
\begin{compactitem}
    \item localized basis set (GTO, NAO, STO, ...)
    \item Resolution-of-Identity (RI) scheme for the
    evaluation of Coulomb integrals
\end{compactitem}
\smallskip

Coulomb vertex can be evaluated via:
\begin{equation*}\label{local-vertex}
  \Gamma_{r\eta}^{q}
  =\sum_{\mu\alpha\gamma}
  {c^\ast}_\alpha^q (\bm{k}_q)c_r^\gamma(\bm{k}_r)
  C_{\gamma\mu}^{\alpha}(\bm{k_q},\bm{k_r})
  \left[\left(V^{\bm{k_q}-\bm{k_r}}\right)^\frac12\right]_\eta^\mu
\end{equation*}
with
\begin{compactitem}
    \item the HF/KS-state vectors $c_{q}^{\alpha}$
    \item the RI-expansion coefficients $C_{\gamma\mu}^{\alpha}$
    \item the auxiliary Coulomb matrix built from the RI-functions
    $P_{\mu}(\bm{r})$:
    $V^{\mu}_{\eta} = \int d\bm{r}\bm{r'} \frac{P_{\eta}(\bm{r}){P^\ast}^{\mu}(\bm{r'})}{|\bm{r}-\bm{r'}|}$
\end{compactitem}

}
\headerbox{Plane waves - VASP}{name=vasp,column=1,below=fhiaims,span=1}{
VASP employs a plane-wave basis set and the projector augmented wave (PAW) method
to approximate the one-electron orbitals with periodic boundary conditions~\cite{blochl1994,kresse1996efficiency}.

The Coulomb vertex can be evaluated via~\cite{hummel2017}:
\begin{equation*}\label{pw-vertex}
\Gamma_{r \bm{G}}^{q}
  =
  \sqrt{\frac{4\pi}{{\bm G}^2}}
  \int{\rm d}{\bm x}\,
  {\rm e}^{-{\rm i}{\bm G}\cdot{\bm r}}\,
  {\psi^\ast}^q({\bm x}) \psi_r({\bm x})\,
\end{equation*}
with
\begin{compactitem}
    \item plane-wave vectors $\bm{G}$
    \item the orbital co-densities ${\psi^\ast}^q({\bm x})\psi_r({\bm x})$
are approximated in the PAW method using Eq.~(2.87) of
Ref.~\cite{harl2008}
    \item
automated basis-set extrapolation for MP2 
    \item
	    automated auxiliary-field truncation~\cite{hummel2017}
    \item
memory-efficient approximate natural orbital algorithm~\cite{gruneis2011}
\end{compactitem}


}


\headerbox{Strong scaling}{name=strongscaling,column=0,below=H2OC3N4,span=1}{
    \begin{center}
		\includegraphics[scale=0.5]{bilder/strong_scaling.pdf}		
	\end{center}
}

\headerbox{Weak scaling}{name=weakscaling,column=1,below=vasp,span=1}{
	\begin{center}
		\includegraphics[scale=0.5]{bilder/weak_scaling.pdf}		
	\end{center}
}

\begin{posterbox}[name=cc4s,column=2,below=einleitung,span=1]{Cc4s}
Coupled cluster for solids (Cc4s):
\begin{compactitem}
\item extended Einstein convention in source, e.g.:
{\small
\begin{lstlisting}
R["abij"] += T["acik"]*T["dblj"]*V["klcd"]
\end{lstlisting}
}
\item
optimal contraction order automatically determined
by tensor-contraction-compiler (\texttt{tcc})
\item
lazy evaluation only of terms needed
\item
interfaced to parallel tensor frameworks,
	currently Cyclops Tensor Framework (CTF)~\cite{solomonik2014} %site...
\item 
standardized input/output for tensors of arbitrary rank
\item efficient CBS and TDL approximations for CCSD~\cite{gruber2018,irmler2021}
\end{compactitem}
\end{posterbox}

\headerbox{Cc4s-Workflow}{name=workflow,column=2,below=cc4s,span=1}{
\begin{center}
	\includegraphics{bilder/cc4s-flowchart}		
\end{center}
}

\headerbox{References}{name=literatur,column=1,below=workflow,span=2}{
\AtNextBibliography{\footnotesize}
\setlength\bibitemsep{0pt}
\printbibliography[heading=none]
}

%\headerbox{Summary and outlook}{name=ergebnis,column=0,below=fhiaims,span=1}{
%\begin{compactitem}
%    \item \texttt{CC4S} and its interfaces are a 
%    valuable tool for materials scientists to obtain reliable \& accurate data
%	\item The range of methods in \texttt{CC4S} is continuously increasing
%	\item The two interfaces pave the way for many other codes - atom-centered and plane wave - to use %\texttt{CC4S}
%\end{compactitem}
%}



\end{poster}
\end{document}
